\chapter{Pendahuluan}

\section{Latar Belakang}
Kemiskinan merupakan fakta sosial global yang selalu menjadi perhatian dari tahun
ke tahun \parencite{megawati2021}. Salah satu metode standar untuk mengidentifikasi
penerima program bantuan sosial adalah \emph{Proxy Means Test} (PMT)
\parencite{etang2022}. \textcite{taufiq2021} menunjukkan bahwa pendekatan
\emph{ensemble} dapat meningkatkan akurasi penargetan.

% (Uraikan konteks, kesenjangan, dan posisi penelitian Anda di sini.)

\section{Identifikasi dan Batasan Masalah}
Bagian ini memuat identifikasi masalah dan batasan ruang lingkup penelitian.

\section{Tujuan Penelitian}
Bagian ini memuat tujuan penelitian yang ingin dicapai.

\section{Manfaat Penelitian}
Bagian ini memuat manfaat teoretis dan praktis dari penelitian.

\section{Sistematika Penulisan}
\ifpeminatansi
Penulisan skripsi ini terdiri atas enam bab. Bab I Pendahuluan menguraikan latar
belakang, identifikasi dan batasan masalah, tujuan, manfaat, serta sistematika
penulisan. Bab II Tinjauan Pustaka memuat landasan teori dan penelitian terkait.
Bab III Metode Penelitian menjelaskan data serta metodologi pengembangan sistem
yang digunakan. Bab IV Analisis dan Perancangan menyajikan analisis kebutuhan dan
rancangan sistem. Bab V Implementasi dan Evaluasi memaparkan implementasi sistem
beserta hasil pengujian dan evaluasinya. Bab VI Kesimpulan dan Saran berisi
kesimpulan penelitian dan saran untuk pengembangan selanjutnya.
\else
Penulisan skripsi ini terdiri atas lima bab. Bab I Pendahuluan menguraikan latar
belakang, identifikasi dan batasan masalah, tujuan, manfaat, serta sistematika
penulisan. Bab II Tinjauan Pustaka memuat landasan teori dan penelitian terkait.
Bab III Metode Penelitian menjelaskan data, tahapan, dan teknik analisis yang
digunakan. Bab IV Hasil dan Pembahasan menyajikan temuan penelitian beserta
pembahasannya. Bab V Kesimpulan dan Saran berisi kesimpulan penelitian dan saran
untuk penelitian selanjutnya.
\fi

% Catatan: paragraf di atas mengikuti penanda \peminatansi (lihat main.tex).
